\documentclass[a4paper,12pt]{article}
\usepackage{amsmath}
\usepackage{graphicx}
\usepackage{multicol}
\usepackage{geometry}
\usepackage{fancyhdr}
\usepackage{tikz}
\usepackage{eso-pic}
\usepackage[spanish]{babel}

\geometry{top=1.5cm, bottom=1.5cm, left=1.5cm, right=1.5cm}

\pagestyle{fancy}
\fancyhf{}
\rhead{Esc. 4-117 Ejército de los Andes}
\lhead{Termodinámica y Máquinas Térmicas}
\cfoot{\thepage}

\begin{document}

\begin{center}
    \Large \textbf{Evaluación de Termodinámica y Máquinas Térmicas} \\
    \normalsize Duración: 80 minutos \\
    Alumno/a: \underline{\hspace{6cm}} \hspace{0.5cm} Fecha: \underline{\hspace{2cm}}
\end{center}

\vspace{0.3cm}

\begin{multicols}{2}

\section*{Problemas}

\begin{enumerate}

\item \textbf{Trabajo externo en sistema cerrado} \\
Un cilindro con un pistón contiene un gas a una presión de $3\,\text{atm}$. Se desea introducir $0{,}5\,\text{m}^3$ de gas en su interior. ¿Cuánto trabajo se necesita realizar para lograrlo?

\vspace{0.3cm}

\item \textbf{Trabajo en almacenamiento de gas} \\
En una instalación de almacenamiento, se desea agregar $1\,\text{m}^3$ de gas a un tanque que ya se encuentra a una presión de $5\,\text{atm}$. Determine el trabajo necesario para efectuar esta tarea.

\vspace{0.3cm}

\item \textbf{Trabajo en sistema de refrigeración} \\
Un sistema de refrigeración emplea un compresor para mantener una presión constante de $2\,\text{atm}$. Si se desea introducir $2\,\text{m}^3$ de gas, ¿cuál es el trabajo que debe realizar el compresor?

\vspace{0.3cm}

\item \textbf{Trabajo efectuado por el gas} \\
Un gas contenido en un cilindro tiene un volumen inicial $v_i = 0{,}025\,\text{m}^3$ y presión $p_i = 2{,}5\,\text{atm}$. Al calentarlo, su volumen aumenta a $v_f = 0{,}175\,\text{m}^3$ manteniéndose constante la presión. Calcule el trabajo efectuado por el gas.

\vspace{0.3cm}

\item \textbf{Trabajo sobre un gas en laboratorio} \\
En un laboratorio, un recipiente sellado contiene gas a una presión de $4\,\text{atm}$. Para un experimento se introduce $0{,}8\,\text{m}^3$ de gas adicional. Determine el trabajo requerido.

\end{enumerate}

\end{multicols}

\end{document}
